% 04-Dataset.tex -- data construction and experimental setup.
%
% AAAI COMPLIANCE: no added packages, and none of the commands aaai2027
% forbids (no \setlength, \vspace, \vskip, \newpage, \tiny, ...). Column
% padding is trimmed with @{} rather than by resetting \tabcolsep.
%
% SCOPE: experimental detail only -- what was run, on what data, scored
% how. No findings, no interpretation; those live in the Results section.

\section{Data}

\subsection{Universe and Source}

All data are daily CRSP records for U.S.\ common equities. From the
mid-- and upper-mid-capitalisation segment we form four disjoint panels,
denoted $V1$--$V4$, of 51 tickers each (204 distinct names in total). Each
panel spans 2004-08-16 to 2024-12-31, or 5{,}130 trading days, and is
stored as one CSV per ticker plus a panel-level date index carrying the
aligned raw price and transaction-cost columns for every name.

Panel membership is fixed for the whole sample: a ticker is either in a
panel for all 5{,}130 days or not in it at all. This removes index-
reconstitution effects from the comparison, at the cost of conditioning
the universe on names that survived to the end of the sample. All methods
trade the identical universe, so the conditioning applies equally to every
arm reported.

\subsection{Features and Prediction Target}

Each model sees seven per-stock features, all computed causally from
information available at or before the formation day:

\begin{itemize}\itemsep2pt
\item \texttt{RET} --- the raw daily return.
\item \texttt{RET\_CS\_IN} --- the cross-sectional rank-normal score of
      \texttt{RET} on the same day, i.e.\ $\Phi^{-1}\!\big((\mathrm{rank}_t(i)-0.5)/N\big)$
      over the $N$ names priced that day.
\item \texttt{BA\_SPREAD} --- the quoted bid--ask spread.
\item \texttt{ILLIQUIDITY} --- Amihud illiquidity: absolute return over
      dollar volume, passed through a 21-day trailing mean and a log.
\item \texttt{REV21\_1} --- short-horizon reversal, the compounded return
      over the 20 rows preceding the formation day, excluding it.
\item \texttt{TURN\_RATIO} --- turnover acceleration: the mean of
      $\log(1+\text{turnover})$ over the last 5 rows minus the same mean
      over the last 63 rows.
\item \texttt{VOL21} --- the sample standard deviation of \texttt{RET}
      over the trailing 21 rows.
\end{itemize}

The last five of these are additionally mapped, per date, to their
cross-sectional rank on $[-1,1]$, so that no feature carries a level that
drifts across the two decades of the sample.

The prediction target is \texttt{RET\_CS}: the next day's return, rank-
normalised across the cross-section of that day by the same $\Phi^{-1}$
transform. Predicting a cross-sectional rank rather than a raw return
makes the target stationary by construction and matches how the
predictions are consumed --- every trading rule in this paper reads only
the ordering of the predicted cross-section.

Realised returns used for scoring are CRSP \texttt{RET}, which is adjusted
for splits and dividends. Prices and per-name transaction costs enter only
through the cost model below.

\section{Experimental Setup}

\subsection{Online Evaluation Protocol}

Every method is evaluated prequentially: on each scored day the model
predicts the next day's cross-section, the realised outcome is then
revealed, and the model updates. No model ever sees a label before
predicting it, and no model is refit on the scored span with hindsight.

ONE-NAS advances on a weekly clock. One generation corresponds to a
\texttt{window\_step} of 5 trading days; each genome is trained on
sequences of length $L=40$ and validated on the 5 most recent windows.
The number of training windows is set per panel so that the clock lands on
the intended calendar span --- 712, 695, 702 and 738 windows for
$V1$--$V4$ on the evaluation span, and 511, 493, 501 and 536 on the
tuning span.

The sample is divided into three disjoint spans, used for three distinct
purposes:

\begin{itemize}\itemsep2pt
\item \textbf{Burn-in}, through 2019-12-31. The population evolves online
      but nothing is scored. A further 50 warm-up generations precede any
      recorded prediction, so the scored population is never a cold start.
\item \textbf{Tuning}, 2016-01-01 to 2019-12-31. Every configuration
      decision --- hyperparameter screen, island count --- is made here.
\item \textbf{Evaluation}, 2020-01-01 to 2024-12-31. Scored once per
      registered configuration.
\end{itemize}

\subsection{Search Configuration and Prediction Rule}

The ONE-NAS settings are listed in Table~\ref{tab:hps} and were selected
on the tuning span. Node types available to the search are simple,
UGRNN, MGU, GRU, $\Delta$-RNN and LSTM recurrent cells.

A single run yields, at every generation, both the elite population of
each island and the single best genome overall, so the prediction rule is
a scoring-time choice rather than a property of the run. Two rules are
reported. The \emph{single champion} rule takes the global-best genome, as
in prior ONE-NAS work. The \emph{ensemble} rule takes one champion per
island --- the lowest-validation-error elite of each --- converts each
member's predicted cross-section to ranks, and averages the ranks. Rank
averaging is used rather than averaging predictions because the members
are separately trained networks whose output scales are not comparable.

\subsection{Baselines}

Four online baselines trade the same books on the same panels under the
same costs. All are single networks; the reported cell is the mean over
independent runs, not the return of an ensemble of those runs.

\begin{itemize}\itemsep2pt
\item \textbf{Online LSTM} --- one LSTM, 8 hidden units, sequence length
      10, Adam at $10^{-3}$, one gradient step per trading day with a
      1000-day replay buffer.
\item \textbf{Online GRU} --- as above with a GRU cell and sequence
      length 5.
\item \textbf{Periodic LSTM} --- the same LSTM architecture refit from
      scratch each month for up to 4000 steps, rather than updated daily.
\item \textbf{Online AR} --- a first-order autoregression fitted by online
      gradient descent at $\eta=10^{-4}$. Deterministic, so it contributes
      one run per panel rather than one per seed.
\end{itemize}

All baselines predict the same cross-sectionally demeaned target as
ONE-NAS.

\subsection{Trading Books}

Predictions are turned into positions by a registered primary
construction, with two further constructions reported as sensitivities.

The \textbf{sleeves} book is the registered primary: Jegadeesh--Titman
overlapping portfolios. Each day a new sleeve goes long the top
$k=10$ names and short the bottom 10 by predicted rank, at
$(\$100/H)/k$ per name per side with holding period $H=10$ days; the
sleeve formed $H$ days earlier is retired. Ten sleeves are therefore live
at once, and costs are charged on the change in \emph{aggregate} per-name
notional, so a name held by several sleeves is netted rather than traded
repeatedly. Because only the ordering of each day's predictions is read,
every metric is invariant to any strictly monotone transform of a model's
outputs.

The \textbf{banded} book enters a name at predicted rank $\le 10$ and
holds it until its rank passes 35, following the buy/hold-spread
construction used to reduce turnover in the transaction-cost literature.
The \textbf{Algorithm 1} book is the conditional long--short rule of the
prior study, which rebalances only when all ten top predictions are
positive and all ten bottom predictions negative.

Books are funded with \$100 of capital and \$200 of gross notional when
fully invested. Transaction costs are charged per name from the CRSP
cost field scaled by the traded notional at the prior close, so
wide-spread names are genuinely more expensive to trade than liquid ones;
no flat basis-point assumption is used. All reported returns are net of
these costs.

\subsection{Metrics}

Two families are reported. \emph{Prediction quality}: the daily
cross-sectional Spearman rank correlation between predicted and realised
returns, averaged over scored days (rank IC). \emph{Economics}: net return
as cumulative daily profit and loss over initial capital, annualised
Sharpe ratio, maximum drawdown, and turnover as mean daily traded notional
divided by gross book size.

Factor-adjusted returns regress each arm's pooled daily book return on the
Fama--French three factors plus momentum, with Newey--West standard errors
at lag 10.

\subsection{Replication and Statistical Comparison}

The unit of replication is one (panel, seed) cell. Unless stated
otherwise, ONE-NAS arms comprise 10 seeds $\times$ 4 panels $=40$ cells.
The baselines were run at higher replication --- 48 seeds for online LSTM,
40 for online GRU and periodic LSTM --- and their own means are reported
at that full replication.

Comparisons between arms are paired within (panel, seed) cells, so that
panel and seed effects cancel, and the reported $t$ is a paired $t$ over
those cells. Where two arms do not share a seed set, the pairing is
restricted to the cells they have in common and $n$ is stated.

\subsection{Compute Environment}

All searches ran on ACCESS Anvil, one 128-core node per run under MPI with
one master and 127 workers. Search cost is measured from job accounting
rather than estimated: a four-panel ONE-NAS run at 40 islands charges 785
core-hours, against 0.03 core-hours for an online-LSTM run and 0.77 for a
periodic-LSTM run. The ONE-NAS figure reflects the charged allocation; the
same configuration completes on approximately 48 core-hours of useful
work, because a majority of wall time is serial evaluation on the master
rank rather than parallel work on the workers.

\subsection{Pre-Registration}

The scored configuration, evaluation span, trading book and metrics were
committed to the repository before the evaluation span was scored, and
each subsequent look at the data was declared in advance with its
decision rule. Configuration changes require a dated amendment written
before the changed configuration is scored. The registered protocol,
the hyperparameter screen with its per-cell outcomes, and every declared
exploratory look are included in the code release.
